\section{A radial Hecke comparison}
\label{v4-arith:radial:section}

The Satake trace recursion also admits a representation on radial functions
on a regular tree.
Its adjacency operator has the density used in
Definition~\ref{v4-arith:def:averaging-naive}. Matching that density
does not identify the operators. The radial tree representation gives an
explicit comparison with the spherical Hecke action and its vector state.
This is the classical spherical transform, with the distance and
symmetric-power normalizations kept separate. The density is the
Kesten--McKay measure. McKay, \emph{The expected eigenvalue distribution
of a large regular graph}, Linear Algebra Appl. 40 (1981),
203--216, Theorem~1.1, gives its graph-theoretic normalization.

Fix a prime $p$ and a rooted $(p+1)$-regular tree $(\mathcal T_p,o)$.
Use counting measure on its vertices and complex coefficients.
Let $S_n$ be the vertices at distance $n$ from $o$, and write
\[
 s_0=1,\qquad s_n=(p+1)p^{n-1}\quad(n\ge1),\qquad
 e_n=s_n^{-1/2}\mathbf{1}_{S_n}.
\]
The radial Hilbert space $\mathcal R_p$ consists of the square-summable
functions constant on each $S_n$. The vectors $e_n$ form an orthonormal
basis. All these vector spaces have cohomological degree zero.
Distance gives a filtration, not a grading preserved by adjacency.

Let $T$ be adjacency, $(Tf)(v)=\sum_{w\sim v}f(w)$.
It is a bounded self-adjoint operator on the vertex Hilbert space,
with norm at most $p+1$ by Cauchy--Schwarz.
Counting edges between consecutive spheres gives
\begin{align}
 Te_0&=\sqrt{p+1}\,e_1,\\
 Te_1&=\sqrt{p+1}\,e_0+\sqrt p\,e_2,\\
 Te_n&=\sqrt p\,(e_{n-1}+e_{n+1})\quad(n\ge2).
 \label{v4-arith:radial:recurrence}
\end{align}
In particular, two steps return to the root in $p+1$ ways.
Thus the polynomial which produces the distance-two sphere is
$T^2-(p+1)I$. The symmetric-power Hecke polynomial is instead
$T^2-pI$. Their difference is the identity operator.

Define monic integral polynomials by
\[
 R_0(x)=1,\quad R_1(x)=x,\quad R_2(x)=x^2-(p+1),\qquad
 R_{n+1}(x)=xR_n(x)-pR_{n-1}(x)\quad(n\ge2).
\]
The preceding recurrence proves
$R_n(T)e_0=\mathbf{1}_{S_n}$. These formulas already give an integral
$\mathbb Z[x]$-module isomorphism
\[
 \mathbb Z[x]\longrightarrow
 \bigoplus_{n\ge0}\mathbb Z\mathbf{1}_{S_n},\qquad
 f\longmapsto f(T)e_0.
\]
Here $x$ acts on the right by $T$. Monicity makes the $R_n$ an integral
basis, so this assertion needs no inversion of $p$.
The normalized Hilbert comparison uses real square roots and complex scalars.

\begin{theorem}[radial spherical transform]
\label{v4-arith:radial:transform}
Put $x(\theta)=2\sqrt p\cos\theta$ for $0\le\theta\le\pi$, and set
\[
 d\nu_p(\theta)=\frac{2p(p+1)\sin^2\theta}
 {\pi(p^2+1-2p\cos2\theta)}\,d\theta,\qquad
 q_n(\theta)=\frac{R_n(x(\theta))}{\sqrt{s_n}}.
\]
The assignment $\mathcal U_pe_n=q_n$ extends uniquely to a unitary
\[
 \mathcal U_p:\mathcal R_p\longrightarrow L^2([0,\pi],\nu_p).
\]
It sends $e_0$ to $1$ and satisfies the commutative square
\[
 \begin{array}{ccc}
 \mathcal R_p &\xrightarrow{T}& \mathcal R_p\\
 \mathcal U_p\downarrow &&\downarrow\mathcal U_p\\
 L^2(\nu_p)&\xrightarrow{M_x}&L^2(\nu_p)
 \end{array}
\]
where $M_x$ is multiplication by $x(\theta)$.
Consequently $\|T|_{\mathcal R_p}\|=2\sqrt p$, and for every bounded
Borel function $f$ on $[-2\sqrt p,2\sqrt p]$,
\[
 \langle e_0,f(T)e_0\rangle
 =\int_0^\pi f(2\sqrt p\cos\theta)\,d\nu_p(\theta).
 \label{v4-arith:radial:state}
\]
\end{theorem}

\begin{proof}
The tridiagonal recurrence makes $e_0$ cyclic for $T$ on $\mathcal R_p$.
We compute its spectral measure to determine the Hilbert norm of the
polynomial map. Remove the first row and column of that tridiagonal
matrix. Its tail has every off-diagonal coefficient equal to $\sqrt p$.
For $\operatorname{Im}z>0$, the first diagonal entry $g(z)$ of its
resolvent $(z-J)^{-1}$ satisfies the Schur-complement identity
\[
 g(z)=\frac1{z-pg(z)}
      =\frac{z-\sqrt{z^2-4p}}{2p}.
\]
The square root is analytic off $[-2\sqrt p,2\sqrt p]$ and is asymptotic
to $z$ at infinity. The asymptotic $g(z)\sim z^{-1}$ selects this root.
The root resolvent therefore equals
\[
 m(z)=\langle e_0,(z-T)^{-1}e_0\rangle
      =\frac1{z-(p+1)g(z)}.
\]
For $|x|<2\sqrt p$, substitution and Stieltjes inversion give
\[
 -\frac1\pi\operatorname{Im}m(x+i0)
 =\frac{(p+1)\sqrt{4p-x^2}}
 {2\pi((p+1)^2-x^2)}.
\]
There are no poles off this interval. A zero of the denominator would
force $g^2=1$ and $z=\pm(p+1)$, whereas the chosen branch has
$g(\pm(p+1))=\pm p^{-1}$. At both endpoints the denominator is
nonzero because $p>1$. Thus $m$ has continuous finite boundary values
on the interval, including its endpoints, and analytic continuation
across its complement. The inversion formula gives only the displayed
absolutely continuous measure. Its total mass is one, since
$m(z)\sim z^{-1}$. Changing variables $x=2\sqrt p\cos\theta$
gives $\nu_p$.

The spectral theorem now gives
\[
 \int R_n(x)R_m(x)\,d\nu_p
 =\langle R_n(T)e_0,R_m(T)e_0\rangle
 =s_n\delta_{nm}.
\]
Hence the $q_n$ are orthonormal. Polynomials are dense in continuous
functions on the compact $x$ interval and thus in its $L^2$ space.
The monic $R_n$ span all polynomials, so the $q_n$ are complete.
This proves that $\mathcal U_p$ is unitary. The recurrence proves
intertwining first on finite sums and then everywhere by boundedness.
The support of the measure is the full closed interval. Multiplication
by $x$ therefore has norm $2\sqrt p$. Bounded Borel functional calculus
proves the state identity.
\end{proof}

\section{Hecke operations and the distinguished state}
\label{v4-arith:radial:hecke}

For the Bruhat--Tits tree of $\mathrm{PGL}_2(\mathbb Q_p)$, vertices are
homothety classes of $\mathbb Z_p$-lattices in $\mathbb Q_p^2$.
Two vertices are adjacent when representatives satisfy
$pL\subset L'\subset L$ with $[L:L']=p$.
The lines in $L/pL$ give exactly $p+1$ neighbors.
With the maximal compact subgroup of Haar volume one, its adjacency
operator is the characteristic-function Hecke operator for
$\operatorname{diag}(p,1)$. Pullback to $\mathrm{GL}_2$ fixes the
central operator $\operatorname{diag}(p,p)$ to be the identity.
Thus the representation has trivial central character.

Define $H_0(x)=1$, $H_1(x)=x$ and
$H_{k+1}(x)=xH_k(x)-pH_{k-1}(x)$ for $k\ge1$.
In the notation of Proposition~\ref{v4-arith:prop:hecke-generating-function},
$H_k(T)$ represents $H_{p^k}$ with $D_p=1$. These are the
complete-symmetric Satake elements in this central character. Their generating series and spectral values are
\[
 \sum_{k\ge0}H_k(x)t^k=(1-xt+pt^2)^{-1},\qquad
 H_k(x(\theta))=p^{k/2}\sum_{j=0}^ke^{i(k-2j)\theta}.
\]
The second formula follows by multiplying the two geometric series
with parameters $\sqrt p\,e^{i\theta}$ and $\sqrt p\,e^{-i\theta}$.
It remains valid at $\theta=0,\pi$, where the parameters coincide.
The transform therefore intertwines every $H_k(T)$ with the stated
multiplication operator, with the same cyclic state.

\begin{proposition}[distance and symmetric-power elements]
\label{v4-arith:radial:distance}
Let $(D_nf)(v)=\sum_{d(v,w)=n}f(w)$ on the vertex Hilbert space.
Then $D_n=R_n(T)$ and
\[
 H_k(T)=\sum_{j=0}^{\lfloor k/2\rfloor}D_{k-2j},\qquad
 \langle e_0,H_k(T)e_0\rangle=
 \begin{cases}1&k\text{ even},\\0&k\text{ odd}.\end{cases}
\]
On $\mathcal R_p$, for $|t|<p^{-1/2}$, the series converges in operator norm and
\[
 \left\langle e_0,(I-tT+pt^2I)^{-1}e_0\right\rangle
 =\frac1{1-t^2}.
\]
\end{proposition}
\begin{proof}
The distance operators are bounded because every row and column has
$s_n$ entries. Counting the first step gives
$TD_1=D_2+(p+1)I$ and
$TD_n=D_{n+1}+pD_{n-1}$ for $n\ge2$. Thus $D_n=R_n(T)$.
With $H_{-1}=H_{-2}=0$, the recurrences give
$R_n=H_n-H_{n-2}$ for every $n\ge0$.
Telescoping proves the sum formula. Since $D_ne_0=\mathbf{1}_{S_n}$,
only $D_0$ contributes to the root state.
The spectral formula bounds $\|H_k(T)|_{\mathcal R_p}\|$
by $(k+1)p^{k/2}$. This proves norm convergence on $\mathcal R_p$
in the stated disk and permits evaluation of its state term by term.
Multiplication of the series by $I-tT+pt^2I$ proves the inverse formula
on that space.
\end{proof}

For example, at $p=2$,
\[
 q_0=1,\qquad q_1=x/\sqrt3,\qquad
 q_2=(x^2-3)/\sqrt6,\qquad H_2=x^2-2.
\]
The tree counts give
$\int x^2d\nu_2=3$ and $\int x^4d\nu_2=15$.
Indeed a four-step return either returns after two steps, giving nine
walks, or first reaches distance two, giving six walks.
Consequently $\int H_2d\nu_2=1$ but $\int R_2d\nu_2=0$,
and $\int R_2^2d\nu_2=6$. These values distinguish the two Hecke
normalizations already at the first nontrivial higher element.

\section{Finite sets of primes and the arithmetic comparison}
\label{v4-arith:radial:finite}

For a finite prime set $F$, form the completed Hilbert tensor products
$\mathcal R_F=\widehat\bigotimes_{p\in F}\mathcal R_p$ and the product
probability measure $\nu_F=\prod_{p\in F}\nu_p$.
The tensor product of the maps $\mathcal U_p$ is the unitary
\[
 \mathcal U_F:\mathcal R_F\longrightarrow
 L^2([0,\pi]^F,\nu_F),\qquad
 \bigotimes_pe_{n_p}\longmapsto\prod_pq_{n_p}(\theta_p).
\]
It intertwines each local Hecke operation, hence every element of
$\bigotimes_{p\in F}\mathbb C[T_p]$, with multiplication by its
Satake polynomial. The tensor root vector maps to $1$, so its vector
state becomes integration against $\nu_F$. The algebraic integral
lattices are the finite tensor products of the sphere-sum lattices above.
Adjoining a prime uses $v\mapsto v\otimes e_0$ on the tree side and
$f\mapsto f\otimes1$ on the spectral side. These isometric maps
commute with $\mathcal U_F$ and preserve the distinguished states.
They give an operation-compatible finite-stage system.

To specify the occupation-phase comparison, let
$P_p=\mathbb C[a_k,b_k\mid k\ge1]$ be the polynomial oscillator space.
For finitely supported nonnegative integers $r_k,s_k$, give its monomial
$\prod_k a_k^{r_k}b_k^{s_k}$ squared norm
$\prod_k k^{r_k+s_k}r_k!s_k!$, with distinct monomials orthogonal.
Complete this positive inner product to a Hilbert space $\mathcal F_p$.
The charge of this monomial is $m=\sum_k(r_k-s_k)$.
Multiplication by $e^{im\theta}$ on each charge space defines a unitary
$V_\theta$ on $\mathcal F_p$. This is the one-prime occupation action in
Definition~\ref{v4-arith:def:averaging-naive}. The strong integral
$A_p=\int_0^\pi V_\theta\,d\nu_p(\theta)$ is bounded with norm at most one.
It acts on a charge-$m$ monomial by
$c_m=\int_0^\pi e^{im\theta}\,d\nu_p(\theta)$.
For $m=1$, symmetry about $\pi/2$ gives
\[
 c_1=i\int_0^\pi\sin\theta\,d\nu_p(\theta)\in i\mathbb R_{>0}.
\]
Thus this averaging operator has a nonreal eigenvalue. It cannot be
unitarily intertwined with the self-adjoint adjacency operator or
with multiplication by its real spectral coordinate.
This obstruction concerns those specified operators. It does not
exclude a different arithmetic comparison.

The radial transform supplies a complex Hilbert Hecke representation,
an integral cyclic module, and a compatible finite-prime state system.
A comparison with the proposed primitive arithmetic trace class
$\mathcal V^{\mathrm{prim}}$ requires constructing
its arithmetic trace carrier and a map preserving the specified Hecke
operations and trace pairing. Tame residue and Gersten operations
require additional maps on common coefficients.
No full Weil functional or global spectral realization is identified
by the positive root state.
